\section{Initial Planning}
\label{ProjMgmtInitialPlanning}

Todo,
Work for Future Roland
He will have to describe the milestones

\subsection{Project plan and work packages}
Todo,
Work for Future Roland
He will have to describe the project plan and the planned workpackages
See at following content to take the needed out of this


 \section*{Master Thesis Outline Roland Roccaro}

This document describes the outline of the Master Thesis of Roland Roccaro. This 
work is condocted with and for stakeholders at Swiss Post DWH LS OPS (Data 
warehouse Logistic Services Operations). The stakeholders define the standard 
reports to be used for all analytical reporting of all mail and parcel related 
workflows such as sorting, delivery etc.

It shows first the previous work, and then goes trough the research question, basis 
project management that will be applied. Following up with the current 
timeline trough out the Master Thesis using milestones and the outline of the 
chapters which are partly based on the project 1 and 2.

\section*{Study structure in regard to project structure.}

The study programme consists of three different project works. These projects do not 
necessarily have to build on each other.

\begin{itemize}
    \item Project 1. September 2025 to January 2026 9 ECTS / 270 hours
    \begin{itemize}
        \item Fundamental research on the description and concepts of data quality.
        \item Analysis and description of the relevant environment at Swiss Post.
        \item Description of the data lineage using an example.
        \item Description of the current manual process.
        \item Description of the desired automated process.
    \end{itemize}

    \item Project 2. February 2026 to August 2026 15 ECTS / 450 hours
    \begin{itemize}
        \item Development and improvement of the report described in Project 1, 
        including the development of several simplified visualisations.
        \item Iterative improvement in collaboration with business stakeholders, 
        including Scrum Review presentations and subsequent improvements.
        \item Testing of monitoring infrastructure to enable the collection of value 
        ranges and support future alerting.
    \end{itemize}

    \item Master Thesis. September 2026 to August 2027 30 ECTS / 900 hours
    \begin{itemize}
        \item This document describes the initial structure of the Master Thesis.
    \end{itemize}
\end{itemize}

 \section*{Research Question}

Following research question is the core of this master thesis:

\begin{itemize}
    \item To what extent can data quality in analytical reporting at DWH LS OPS be 
    automatically validated at the reporting layer?
\end{itemize}

 \section*{Hypothesis}

Based on the research question the following hypothesis are defined.

\begin{itemize}
    \item Relevant data quality errors in analytical reporting can be systematically 
    identified, classified, and documented by the responsible stakeholders.
    
    \item An automated data quality validation approach can reliably detect the 
    documented error classes at the analytical reporting layer.
\end{itemize}


 \section*{Research Method}

The research is carried out in phases.

\begin{itemize}
    \item The project starts with an initial research phase, during which a baseline and 
    direction is established.
    
    \item The results are documented as baseline.
    
    \item During the course of the project when reaching relevant phases or milestones, 
    additional research is performed to validate results.
    
    \item Furthermore, after completing the first iteration, further research confirming 
    topics for a second iteration.
\end{itemize}

 \section*{Project Management Method}

The overall project management methods to be used during the master thesis. This 
does not mention all project management methods such as risk management.

\begin{itemize}

    \item Overall Structure Based on Milestones
    \begin{itemize}
        \item The milestones are used to structure the different stages of the research 
        project, from the scientific foundation through implementation to the 
        subsequent evaluation.
        
        \item The milestones are defined at the beginning of the thesis.
        
        \item They divide the work into blocks that provide an overall framework 
        without predetermining the specific topics or approaches.
        
        \item They also provide a basis for subsequently analysing and documenting 
        the effort required for the individual work blocks.
    \end{itemize}

    \item Weekly Meetings, Jour Fixe of Student and advisor
    \begin{itemize}
        \item Ken as advisor and Roland as Student meet weekly for 1 hour
        
        \item Topic of the meeting is:
        \begin{itemize}
            \item Check-in, what is the project, what deliverables are done.
            \item Which Deliverable am I working on?
            \item State of the Deliverable
            \item Next Steps
            \item Other
        \end{itemize}
        
        \item Uppon reaching milestones or other specific dates, the topics can vary.
        
        \item The expert has been invited to the meeting as optional.
    \end{itemize}

    \item Agile Project Management using Scrum.
    \begin{itemize}
        \item Two-week sprints. Every uneven week number sprints start and end.
        
        \item The sprints are synchronised with those of the Swiss Post DWH LS OPS 
        teams.
        
        \item Sprint planning is conducted together with the relevant Swiss Post 
        stakeholders.
        
        \item The planning is reviewed and confirmed during the biweekly Jour Fixe 
        meetings.
        
        \item Scrum Reviews take place every four weeks within DWH LS OPS. These 
        meetings can also be used to present and discuss the progress of the 
        Master Thesis.
        
        \item f. Scrum.
        
        \item Retrospectives take place every two weeks within the DWH LS OPS 
        teams.
        
        \item g. Within the academic team, retrospectives are planned upon 
        completion of each milestone.
    \end{itemize}

\end{itemize}

 \section*{Phases / Milestones}

The following phase are defining blocks which can be closed after completing each 
and move on. The conclusion and start of each phase are used to review and adapt the
current state.

The milestones are events during the course of the phases. Each milestone will be 
documented if it has been reached in time. Furthermore, any adaptions required to 
the milestone plans have to be documented.

\begin{itemize}

    \item Phase 1, Research
    \begin{itemize}
        \item Refine the problem statement.
        \item Define and narrow the scope.
        \item Define the use case.
        \begin{itemize}
            \item Milestone 1, use case defined.
        \end{itemize}
        \item Conduct the literature review.
        \item Define the basic outlines for P2 to P3 and specify possible alternatives.
        \item Closure conditions:
        \begin{itemize}
            \item Research based on problem statement initially closed.
            \item Basis for P2 and P3 defined.
            \item Methodology for thesis defined.
            \item Review of baseline document parts from project 1 and 2 
            performed.
        \end{itemize}
    \end{itemize}

    \item Phase 2, Define the Baseline with Stakeholders
    \begin{itemize}
        \item Adapt M2 based on the findings from P1.
        \item Document the current stakeholder knowledge regarding error detection.
        \begin{itemize}
            \item Milestone 2, Stakeholder meetings held.
        \end{itemize}
        \item Attempt to systematically identify and document relevant errors.
        \item Define how the results will be evaluated against future iterations.
        \begin{itemize}
            \item Milestone 3, Error documentation by stakeholders started.
        \end{itemize}
        \item Closure Conditions
        \begin{itemize}
            \item Stakeholder knowledge documented
            \item Stakeholders have means to document errors.
            \item Defined test window during which manual and automatic errors 
            are compared.
            \item Error documentation setup so that it is comparable to future 
            Phase 3 implementation.
        \end{itemize}
    \end{itemize}

    \item Phase 3, First Iteration
    \begin{itemize}
        \item Review the results of P2 and adapt P3 accordingly.
        \item Define a minimal outline of the P3 objective while keeping the specific 
        approach open to findings from the preceding research.
        \begin{itemize}
            \item Milestone 4, Development environment set, initial setup of 
            tools
        \end{itemize}
        \item Develop the first/basic implementation.
        \begin{itemize}
            \item Milestone 5, MVP has been setup.
        \end{itemize}
        \item Summarise the results.
        \begin{itemize}
            \item Milestone 6, MVP presented to Post stakeholders.
        \end{itemize}
        \item Closure conditions
        \begin{itemize}
            \item The development environment is setup.
            \item Permissions required to work have been setup
            \item The MVP is setup and running.
            \item Data for test window retrieved.
            \item The Post stakeholders received presentation of current state.
            \item Defined test slot and registered test results
        \end{itemize}
    \end{itemize}

    \item Phase 4, Second Iteration
    \begin{itemize}
        \item Review the achievement of the objectives from P3 and, if necessary, 
        pursue an alternative approach.
        \begin{itemize}
            \item Milestone 7, Outline of second iteration defined.
        \end{itemize}
        \item Define a minimal outline of the phase objective while keeping the 
        specific approach open to findings from the preceding research.
        \item Develop an adapted implementation.
        \item Summarise the results.
        \item Closure conditions
        \begin{itemize}
            \item Second iteration outline defined.
            \item Data for test window retrieved.
        \end{itemize}
    \end{itemize}

    \item Phase 5, Comparative Evaluation
    \begin{itemize}
        \item Compare the results of all investigated approaches.
        \begin{itemize}
            \item Milestone 8, Data of different approaches compared.
        \end{itemize}
        \item Identify factors that may limit a direct or fair comparison, including 
        specific differences between the approaches.
    \end{itemize}

    \item Phase 6, Analysis of Goal Achievement and Thesis Finalisation
    \begin{itemize}
        \item Describe which objectives were achieved and which were not.
        \item Finalise the thesis document.
        \begin{itemize}
            \item Milestone 9, Thesis document completed.
        \end{itemize}
        \item Finalise the required deliverables:
        \begin{itemize}
            \item Video
            \item Poster
            \item Presentation and Defence
        \end{itemize}
        \item Milestone 10, Presentation ready
        \item Thesis document delivered
        \item Presentation held
    \end{itemize}

\end{itemize}

 \section*{Milestones Timeline}

The milestone timelines are setup as follows before project start.

\begin{itemize}
    \item Milestone 1, use case defined: 18.10.2026
    \item Milestone 2, Stakeholder meetings held: 15.11.2026
    \item Milestone 3, Error documentation by stakeholders started: 29.11.2026
    \item Milestone 4, Development environment set, initial setup of tools: 13.12.2026
    \item Milestone 5, MVP has been setup: 21.03.2027
    \item Milestone 6, MVP presented to Post stakeholders: 04.04.2027
    \item Milestone 7, Outline of second iteration defined: 18.04.2027
    \item Milestone 8, Data of different approaches compared: 16.05.2027
    \item Milestone 9, Thesis document completed: 11.07.2027
    \item Milestone 10, Presentation ready: 25.07.2027
\end{itemize}